\answer{OA03}
记 $A=\{X+Y>1\}$。原来的联合密度在单位正方形上恒为 $1$，而 $\Prob(A)=1/2$，故条件密度在上三角形内恒为 $2$。给定 $X=x$ 与 $A$ 后，$Y$ 在 $(1-x,1)$ 上均匀分布，于是
\[
 \E[Y\mid X=x,A]=1-\frac x2,
 \qquad 0<x<1.
\]
条件均值已经是仿射函数。令 $u=Y-(1-X/2)$，则 $\E[u\mid X,A]=0$，从而 $u$ 与常数及 $X$ 正交。因此最小二乘回归的截距为 $1$、斜率为
\[
 \boxed{-\frac12}.
\]
也可用矩核对。条件密度 $f_{X\mid A}(x)=2x$ 给出
\[
 \E[X\mid A]=\frac23,\quad
 \E[X^2\mid A]=\frac12,\quad
 \Var(X\mid A)=\frac1{18}.
\]
由条件均值公式，$\Cov(X,Y\mid A)=-\Var(X\mid A)/2=-1/36$，协方差除以方差仍得到 $-1/2$。独立性属于筛选前的总体，不能直接用于条件总体。
